% =====================================================================
%  config/langues.tex
%  Gestion multilingue avec polyglossia (adapté à XeLaTeX).
%
%  Langue principale : FRANÇAIS (typographie française : espaces
%  insécables avant : ; ! ?, guillemets « … », etc.).
%
%  Langues secondaires pré-configurées, avec leur écriture :
%    - arabe    (arabe,   RTL — vrai rendu droite-à-gauche via bidi)
%    - russe    (cyrillique)
%    - chinois  (hans/CJK)
%    - japonais (CJK)
%    - grec     (grec moderne/ancien)
%    - anglais  (pour l'abstract)
%
%  IMPORTANT : polyglossia doit être chargé APRÈS fontspec (polices.tex)
%  et les familles \arabicfont, \russianfont, etc. doivent déjà exister.
%
%  ARABE + TABLEAUX : l'option « arabic » charge le paquet « bidi », qui
%  assure un vrai algorithme de paragraphe RTL (alignement à droite,
%  césure, ordre des lignes). Sous TeX Live 2026 (dont Overleaf), bidi
%  provoque une erreur « \UseMathForPositioningText » dans les tableaux ;
%  le correctif est appliqué dans config/preambule.tex (shim), chargé
%  AVANT bidi. Voir ce fichier pour les détails.
% =====================================================================

\usepackage{polyglossia}

% --- Langue principale -----------------------------------------------
\setmainlanguage{french}

% --- Langues secondaires ---------------------------------------------
%  polyglossia associe automatiquement \arabicfont, \russianfont,
%  \chinesefont, \japanesefont, \greekfont (définis dans polices.tex).
\setotherlanguage{english} % pour l'abstract en anglais
\setotherlanguage{arabic}
\setotherlanguage{russian}
\setotherlanguage[variant=simplified]{chinese}
\setotherlanguage{japanese}
\setotherlanguage{greek}

% ---------------------------------------------------------------------
%  Commandes de confort pour insérer un court fragment dans une écriture
%  donnée, sans avoir à ouvrir un environnement complet.
%
%  Exemples d'utilisation :
%     \ar{النص العربي}     % arabe (RTL géré par polyglossia/bidi)
%     \ru{Русский текст}   % russe
%     \zh{中文文本}         % chinois
%     \ja{日本語のテキスト}  % japonais
%     \gr{Ἑλληνικά}        % grec
% ---------------------------------------------------------------------
\newcommand{\ar}[1]{\textarabic{#1}}
\newcommand{\ru}[1]{\textrussian{#1}}
\newcommand{\zh}[1]{\textchinese{#1}}
\newcommand{\ja}[1]{\textjapanese{#1}}
\newcommand{\gr}[1]{\textgreek{#1}}

% ---------------------------------------------------------------------
%  Passages longs
%  ---------------------------------------------------------------------
%  Utilisez les environnements fournis par polyglossia, qui gèrent la
%  césure et la direction du texte :
%     \begin{russian} ... \end{russian}
%     \begin{chinese} ... \end{chinese}   etc.
%
%  Pour un paragraphe ARABE, utilisez l'environnement « Arabic » avec une
%  MAJUSCULE (pour éviter le conflit avec la commande \arabic de LaTeX) :
%     \begin{Arabic} ... \end{Arabic}
%  Il assure un véritable rendu droite-à-gauche du bloc (alignement à
%  droite, retours à la ligne et césure corrects).
% ---------------------------------------------------------------------
